\documentclass[11pt,a4paper]{article}

\usepackage[spanish,es-noshorthands]{babel}
\usepackage[T1]{fontenc}
\usepackage[utf8]{inputenc}
\usepackage[a4paper,margin=2.4cm]{geometry}
\usepackage{amsmath,amssymb}
\usepackage{listings}
\usepackage{float}
\usepackage{array}
\usepackage{booktabs}
\usepackage{tikz}
\usetikzlibrary{automata,positioning}
\usepackage[colorlinks=true,linkcolor=blue,urlcolor=blue]{hyperref}

\lstset{
  basicstyle=\ttfamily\small,
  frame=single,
  breaklines=true,
  columns=fullflexible,
  keepspaces=true,
  showstringspaces=false
}

\title{\textbf{Ejercicio Práctico 4}\\[2mm]
\large Sistema de Transición de Estados en NuSMV: Búsqueda de Rutas}
\author{Proyecto: Sistema de Planificación de Viajes en Transporte Público (aqpTracker)}
\date{\today}

\begin{document}

\maketitle

\section{Objetivo}

Modelar como sistema de transición de estados el ciclo de una solicitud de búsqueda de rutas de transporte público, identificando estados, variables, condiciones de transición y actualizaciones de valores, y verificando un requerimiento mediante los operadores temporales \texttt{G} y \texttt{F} de NuSMV.

\section{¿Qué es un sistema de transición de estados?}

Un sistema de transición de estados (STS, por sus siglas en inglés) es una estructura formal:
\[
M = (S, S_0, \rightarrow, AP, L)
\]
donde \(S\) es un conjunto finito de estados, \(S_0 \subseteq S\) son los estados iniciales, \(\rightarrow \subseteq S \times S\) es la relación de transición, \(AP\) es un conjunto de proposiciones atómicas y \(L : S \rightarrow 2^{AP}\) indica qué proposiciones son verdaderas en cada estado.

En NuSMV el sistema se describe con variables de dominio finito, asignaciones iniciales (\texttt{init}) y asignaciones de siguiente estado (\texttt{next}). Las propiedades se escriben en LTL:
\begin{itemize}
  \item \texttt{G} \(\varphi\) (``globalmente \(\varphi\)''): \(\varphi\) se cumple en todos los estados de todos los caminos. Se usa para propiedades de seguridad o invariantes.
  \item \texttt{F} \(\varphi\) (``eventualmente \(\varphi\)''): \(\varphi\) se cumple en algún estado futuro de todos los caminos. Se usa para propiedades de vivacidad o alcanzabilidad.
\end{itemize}

\subsection{Contexto metodológico}

La traducción de modelos de estado a verificadores simbólicos se ha aplicado en diversos dominios. Kochaleema y Santhosh Kumar (2022) traducen Statecharts UML a sistemas de transición etiquetados y verifican propiedades CTL/LTL; Hlaoui y Ayari (2025) compilan modelos de proceso correctos por construcción a NuSMV; y Xavier et al. (2024) sintetizan máquinas de estado desde registros de eventos y las verifican con NuSMV. Este trabajo adopta el mismo enfoque: un modelo abstracto de estados, propiedades LTL y verificación con NuSMV.

\section{Descripción del modelo}

Se modela el ciclo de una solicitud de búsqueda de rutas. El sistema parte en \texttt{inicio}, esperando una solicitud. Si la solicitud es válida, pasa a \texttt{buscando}; si es inválida, pasa a \texttt{solicitud\_invalida}. Cuando la búsqueda termina, el sistema pasa a \texttt{con\_ruta} si se encontró al menos una ruta, o a \texttt{sin\_ruta} si no se encontró ninguna. Los estados \texttt{con\_ruta}, \texttt{sin\_ruta} y \texttt{solicitud\_invalida} son transitorios: después de mostrar el resultado, el sistema vuelve a \texttt{inicio}.

La variable \texttt{solicitud} representa la entrada del usuario y se actualiza de forma no determinista, modelando que en cada paso puede llegar una solicitud válida, una inválida o ninguna. La variable \texttt{rutas\_encontradas} se asigna al iniciar una búsqueda válida y determina el resultado. La variable \texttt{consultas} cuenta las solicitudes válidas procesadas.

El modelo es una abstracción del flujo de los endpoints de búsqueda directa y combinada: el backend actual es stateless (responde a una petición POST), pero la interfaz de usuario sí atraviesa estos estados. El modelo sirve como especificación del comportamiento esperado y como base para una futura máquina de estados en el frontend o en un servicio de búsqueda. La distinción entre búsqueda directa y combinada se deja al algoritmo (VDM++); la máquina de estados modela el ciclo común de la solicitud.

\section{Estados, variables y estado inicial}

\subsection{Estados}

\begin{center}
\begin{tabular}{|l|p{9cm}|}
\hline
\textbf{Estado} & \textbf{Significado} \\
\hline
\texttt{inicio} & Esperando una solicitud de búsqueda \\
\hline
\texttt{buscando} & Solicitud válida en proceso; abstrae búsqueda directa y combinada \\
\hline
\texttt{con\_ruta} & La búsqueda encontró al menos una ruta \\
\hline
\texttt{sin\_ruta} & La búsqueda no encontró rutas \\
\hline
\texttt{solicitud\_invalida} & La solicitud es inválida \\
\hline
\end{tabular}
\end{center}

\subsection{Variables}

\begin{center}
\small
\begin{tabular}{|l|p{5.8cm}|p{6.2cm}|}
\hline
\textbf{Variable} & \textbf{Dominio} & \textbf{Descripción} \\
\hline
\texttt{estado} & \texttt{\{inicio, buscando, con\_ruta, sin\_ruta, solicitud\_invalida\}} & Estado actual del proceso \\
\hline
\texttt{solicitud} & \texttt{\{ninguna, valida, invalida\}} & Entrada del usuario \\
\hline
\texttt{rutas\_encontradas} & \texttt{0..10} & Número de rutas devueltas por la búsqueda \\
\hline
\texttt{consultas} & \texttt{0..10} & Número de solicitudes válidas procesadas \\
\hline
\end{tabular}
\end{center}

Estado inicial: \texttt{estado = inicio}, \texttt{solicitud = ninguna}, \texttt{rutas\_encontradas = 0}, \texttt{consultas = 0}.

\subsection{Diagrama de estados}

\begin{center}
\begin{tikzpicture}[->,>=stealth,shorten >=1pt,auto,node distance=3.2cm,semithick]
  \node[state,initial] (inicio) {\texttt{inicio}};
  \node[state] (buscando) [right=of inicio] {\texttt{buscando}};
  \node[state] (conruta) [right=of buscando] {\texttt{con\_ruta}};
  \node[state] (sinruta) [below=of buscando] {\texttt{sin\_ruta}};
  \node[state] (invalida) [below=of inicio] {\texttt{solicitud\_invalida}};
  \path
    (inicio) edge node[above]{\texttt{solicitud} = \texttt{valida}} (buscando)
    (inicio) edge node[left]{\texttt{solicitud} = \texttt{invalida}} (invalida)
    (buscando) edge node[above]{\texttt{rutas\_encontradas} $>$ 0} (conruta)
    (buscando) edge node[right]{\texttt{rutas\_encontradas} = 0} (sinruta)
    (conruta) edge[bend left=25] node[above]{resultado} (inicio)
    (sinruta) edge[bend right=25] node[below]{resultado} (inicio)
    (invalida) edge[bend right=25] node[left]{reinicio} (inicio);
\end{tikzpicture}
\end{center}

\section{Condiciones de transición y actualizaciones}

\begin{center}
\begin{tabular}{|p{2.4cm}|p{4.6cm}|p{2.6cm}|p{4.6cm}|}
\hline
\textbf{Estado actual} & \textbf{Condición} & \textbf{Estado siguiente} & \textbf{Actualización} \\
\hline
\texttt{inicio} & \texttt{solicitud = valida} & \texttt{buscando} & \texttt{consultas := consultas + 1}; se asigna \texttt{rutas\_encontradas} \\
\hline
\texttt{inicio} & \texttt{solicitud = invalida} & \texttt{solicitud\_invalida} & sin cambios en contadores \\
\hline
\texttt{buscando} & \texttt{rutas\_encontradas > 0} & \texttt{con\_ruta} & sin cambios \\
\hline
\texttt{buscando} & \texttt{rutas\_encontradas = 0} & \texttt{sin\_ruta} & sin cambios \\
\hline
\texttt{con\_ruta} & --- & \texttt{inicio} & sin cambios \\
\hline
\texttt{sin\_ruta} & --- & \texttt{inicio} & sin cambios \\
\hline
\texttt{solicitud\_invalida} & --- & \texttt{inicio} & sin cambios \\
\hline
\end{tabular}
\end{center}

Actualizaciones de las variables en cada transición:

\begin{center}
\begin{tabular}{|l|p{11cm}|}
\hline
\textbf{Variable} & \textbf{Actualización} \\
\hline
\texttt{solicitud} & Si \texttt{estado = inicio}, se actualiza de forma no determinista a \texttt{ninguna}, \texttt{valida} o \texttt{invalida}; en otro caso, permanece. \\
\hline
\texttt{rutas\_encontradas} & Si \texttt{estado = inicio} y \texttt{solicitud = valida}, se asigna un valor no determinista en \texttt{0..3}; en otro caso, permanece. \\
\hline
\texttt{consultas} & Si \texttt{estado = inicio} y \texttt{solicitud = valida} y \texttt{consultas < 10}, se incrementa en 1; en otro caso, permanece. \\
\hline
\end{tabular}
\end{center}

La asignación de \texttt{rutas\_encontradas} al iniciar la búsqueda permite que el estado \texttt{buscando} clasifique el resultado en el paso siguiente. La variable \texttt{consultas} solo aumenta con solicitudes válidas, de modo que una solicitud inválida no cuenta como consulta.

\section{Código NuSMV}

El código completo del modelo se presenta en la Figura~\ref{fig:codigo-rutas}.

\begin{figure}[H]
\centering
\begin{lstlisting}[basicstyle=\ttfamily\footnotesize]
MODULE main
VAR
  estado            : {inicio, buscando, con_ruta, sin_ruta, solicitud_invalida};
  solicitud         : {ninguna, valida, invalida};
  rutas_encontradas : 0..10;
  consultas         : 0..10;

ASSIGN
  init(estado)            := inicio;
  init(solicitud)         := ninguna;
  init(rutas_encontradas) := 0;
  init(consultas)         := 0;

  next(estado) := case
    estado = inicio & solicitud = valida      : buscando;
    estado = inicio & solicitud = invalida    : solicitud_invalida;
    estado = buscando & rutas_encontradas > 0 : con_ruta;
    estado = buscando & rutas_encontradas = 0 : sin_ruta;
    estado = con_ruta                         : inicio;
    estado = sin_ruta                         : inicio;
    estado = solicitud_invalida                            : inicio;
    TRUE                                      : estado;
  esac;

  next(solicitud) := case
    estado = inicio : {ninguna, valida, invalida};
    TRUE            : solicitud;
  esac;

  next(rutas_encontradas) := case
    estado = inicio & solicitud = valida : {0, 1, 2, 3};
    TRUE                                 : rutas_encontradas;
  esac;

  next(consultas) := case
    estado = inicio & solicitud = valida & consultas < 10 : consultas + 1;
    TRUE                                                   : consultas;
  esac;

DEFINE
  resultado := (estado = con_ruta | estado = sin_ruta);

LTLSPEC G (estado = con_ruta -> rutas_encontradas > 0);
LTLSPEC G (estado = sin_ruta -> rutas_encontradas = 0);
LTLSPEC G (estado = buscando -> F resultado);
LTLSPEC G (estado = inicio & solicitud = valida -> F resultado);
LTLSPEC G (estado = con_ruta -> X estado = inicio);
LTLSPEC G (estado = sin_ruta -> X estado = inicio);
LTLSPEC G (estado = solicitud_invalida -> X estado = inicio);
LTLSPEC G (resultado -> consultas > 0);
\end{lstlisting}
\caption{Código NuSMV del sistema de transición de estados para la búsqueda de rutas.}
\label{fig:codigo-rutas}
\end{figure}

La definición \texttt{resultado} es un alias que abrevia \texttt{estado = con\_ruta | estado = sin\_ruta}; se usa para simplificar la lectura de las propiedades.

\section{Requerimiento verificado con G y F}

\subsection{Requerimiento principal}

\textbf{R3:} ``Toda solicitud válida eventualmente alcanza un resultado: \texttt{con\_ruta} o \texttt{sin\_ruta}.''

Fórmula LTL:

\begin{lstlisting}
LTLSPEC G (estado = inicio & solicitud = valida -> F resultado);
\end{lstlisting}

La fórmula combina \texttt{G} y \texttt{F}: \texttt{G} exige que la implicación se cumpla en todos los estados alcanzables, y \texttt{F} exige que, cuando llega una solicitud válida, exista un estado futuro en el que la búsqueda haya terminado con o sin rutas. La propiedad se cumple porque una solicitud válida en \texttt{inicio} conduce a \texttt{buscando} y, en el paso siguiente, \texttt{rutas\_encontradas} determina \texttt{con\_ruta} o \texttt{sin\_ruta}. Además, la condición antecedente es alcanzable (\texttt{SPEC EF (estado = inicio \& solicitud = valida)} devuelve \texttt{true}), por lo que la propiedad no es vacua.

\subsection{Requerimientos complementarios}

\begin{center}
\begin{tabular}{|l|p{5.4cm}|p{5.4cm}|l|}
\hline
\textbf{ID} & \textbf{Requerimiento} & \textbf{Fórmula LTL} & \textbf{Tipo} \\
\hline
R1 & No se declara éxito sin rutas & \texttt{G (estado = con\_ruta -> rutas\_encontradas > 0)} & Seguridad \\
\hline
R2 & No se declara fracaso con rutas & \texttt{G (estado = sin\_ruta -> rutas\_encontradas = 0)} & Seguridad \\
\hline
R4 & Toda búsqueda en curso termina & \texttt{G (estado = buscando -> F resultado)} & Vivacidad \\
\hline
R5 & Los resultados son transitorios & \texttt{G (estado = con\_ruta -> X estado = inicio)} y \texttt{G (estado = sin\_ruta -> X estado = inicio)} & Seguridad \\
\hline
R6 & Una solicitud inválida vuelve al inicio & \texttt{G (estado = solicitud\_invalida -> X estado = inicio)} & Seguridad \\
\hline
R7 & Un resultado implica al menos una consulta & \texttt{G (resultado -> consultas > 0)} & Seguridad \\
\hline
\end{tabular}
\end{center}

\section{Ejecución y resultados}

Ejecutar:

\begin{lstlisting}
NuSMV rutas.smv
\end{lstlisting}

Salida esperada:

\begin{lstlisting}
-- specification  G (estado = con_ruta -> rutas_encontradas > 0)  is true
-- specification  G (estado = sin_ruta -> rutas_encontradas = 0)  is true
-- specification  G (estado = buscando ->  F resultado)  is true
-- specification  G ((estado = inicio & solicitud = valida) ->  F resultado)  is true
-- specification  G (estado = con_ruta ->  X estado = inicio)  is true
-- specification  G (estado = sin_ruta ->  X estado = inicio)  is true
-- specification  G (estado = solicitud_invalida ->  X estado = inicio)  is true
-- specification  G (resultado -> consultas > 0)  is true
\end{lstlisting}

Todas las propiedades se verifican. Esto significa que, para todas las ejecuciones posibles del modelo, se cumplen las invariantes y las propiedades de vivacidad indicadas.

\subsection{Prueba adicional: una propiedad que no se cumple}

Si se agrega la fórmula:

\begin{lstlisting}
LTLSPEC G (estado = inicio -> F estado = con_ruta);
\end{lstlisting}

NuSMV responde \texttt{is false} y devuelve un contraejemplo: una ejecución en la que \texttt{solicitud} permanece en \texttt{ninguna} y el sistema se queda en \texttt{inicio} indefinidamente. En esa ejecución nunca se alcanza \texttt{con\_ruta}. Esto no es un fallo del modelo: refleja que no toda solicitud produce una ruta; puede no haber solicitud, puede ser inválida o puede no haber rutas disponibles. El requerimiento correcto es R3 (alcanzar un resultado), no ``toda solicitud encuentra una ruta''.

El modelo con esta propiedad adicional y su salida se conservan en \texttt{rutas\_propiedad\_falsa.smv} y \texttt{salida\_rutas\_falsa\_nusmv.txt}.

Traza resumida del contraejemplo:

\begin{lstlisting}
Estado 1.1: estado = inicio, solicitud = ninguna, rutas_encontradas = 0, consultas = 0
Estado 1.2: estado = inicio, solicitud = ninguna, rutas_encontradas = 0, consultas = 0
\end{lstlisting}

\section{Relación con el proyecto}

\begin{center}
\begin{tabular}{|p{4cm}|p{10cm}|}
\hline
\textbf{Elemento NuSMV} & \textbf{Elemento del proyecto} \\
\hline
\texttt{solicitud} & Petición POST a \texttt{buscar\_rutas\_view} o \texttt{buscar\_rutas\_combinadas\_view} \\
\hline
\texttt{rutas\_encontradas} & Longitud de \texttt{rutas\_directas} o \texttt{rutas\_combinadas} en la respuesta \\
\hline
\texttt{consultas} & Contador de búsquedas realizadas \\
\hline
\texttt{estado} & Estados de la interfaz: espera, cargando, resultados, sin resultados, solicitud inválida \\
\hline
\texttt{con\_ruta} & RF1/RF2 con al menos una ruta devuelta \\
\hline
\texttt{sin\_ruta} & RF1/RF2 sin rutas que cumplan las restricciones \\
\hline
\end{tabular}
\end{center}

Las variables \texttt{consultas} y \texttt{estado} son de seguimiento en el modelo; el backend actual no las persiste. El backend actual es stateless: cada petición se procesa y se responde. El modelo no afirma que el backend tenga una máquina de estados persistente; especifica el comportamiento observable de una solicitud y puede servir para implementar una máquina de estados en el frontend o en un futuro servicio de búsqueda. La propiedad R1 garantiza que la interfaz no muestre ``ruta encontrada'' si el resultado está vacío; R2 garantiza que no muestre ``sin ruta'' si hay resultados.

\section{Justificación del diseño}

Esta sección explica por qué el modelo tiene estos estados, variables, condiciones y asignaciones, que es el punto central de la presentación solicitada.

\subsection{¿Por qué estos estados?}

\begin{itemize}
  \item \texttt{inicio}: representa el punto de entrada y el estado de reposo del sistema; sin él no se distingue ``esperando solicitud'' de ``procesando''.
  \item \texttt{buscando}: separa la recepción de la solicitud de la obtención del resultado; modela el tiempo de espera y permite verificar que toda búsqueda termina.
  \item \texttt{con\_ruta} y \texttt{sin\_ruta}: son los dos resultados posibles de RF1/RF2. Distinguirlos evita mostrar ``sin ruta'' cuando sí hay resultados y viceversa.
  \item \texttt{solicitud\_invalida}: separa una solicitud inválida de una búsqueda sin resultados. Sin este estado, una solicitud inválida se confundiría con ``no hay rutas''.
  \item Los estados de resultado son transitorios: la interfaz muestra el resultado y vuelve a \texttt{inicio}, por lo que no se necesitan estados persistentes.
  \item No se separan \texttt{buscando\_directas} y \texttt{buscando\_combinadas} porque la diferencia es algorítmica (VDM++), no del ciclo de vida de la solicitud.
\end{itemize}

\subsection{¿Por qué estas variables?}

\begin{itemize}
  \item \texttt{estado}: variable de control que determina la transición y la respuesta de la interfaz.
  \item \texttt{solicitud}: modela la entrada del usuario; es la condición que dispara la validación y la búsqueda.
  \item \texttt{rutas\_encontradas}: es el dato que decide si el resultado es \texttt{con\_ruta} o \texttt{sin\_ruta}.
  \item \texttt{consultas}: registra cuántas solicitudes válidas se procesaron; sirve para trazabilidad y para verificar que una solicitud inválida no cuenta como consulta.
  \item No se modelan coordenadas, radio ni transbordos porque son parámetros del algoritmo de búsqueda, no del ciclo de vida de la solicitud; esos elementos están en el modelo VDM++.
\end{itemize}

\subsection{¿Por qué estas condiciones?}

\begin{itemize}
  \item \texttt{solicitud = valida} / \texttt{solicitud = invalida}: separan la validación de la búsqueda; una solicitud inválida no debe llegar al buscador.
  \item \texttt{rutas\_encontradas > 0} / \texttt{rutas\_encontradas = 0}: clasifican el resultado de la búsqueda en éxito o ausencia de rutas.
  \item La condición de retorno a \texttt{inicio} no depende de una variable adicional: una vez mostrado el resultado, el ciclo vuelve a empezar.
\end{itemize}

\subsection{¿Por qué estas asignaciones?}

\begin{itemize}
  \item \texttt{rutas\_encontradas} se asigna al iniciar la búsqueda: así el estado \texttt{buscando} puede clasificar el resultado en el paso siguiente y la asignación es no determinista para cubrir todos los resultados posibles.
  \item \texttt{consultas} solo incrementa con solicitudes válidas: evita contar intentos inválidos y mantiene la coherencia con \texttt{resultado -> consultas > 0}.
  \item \texttt{solicitud} solo cambia en \texttt{inicio}: evita que la entrada cambie mientras la búsqueda está en curso.
  \item Los estados de resultado vuelven a \texttt{inicio}: garantizan que la interfaz no quede bloqueada y que el sistema pueda recibir otra solicitud.
\end{itemize}

\subsection{Conclusión de la justificación}

El diseño mantiene el modelo mínimo necesario para representar el ciclo de una solicitud, separa validación, búsqueda y resultado, y deja los detalles algorítmicos (rutas directas/combinadas, radios, transbordos) en VDM++. Esto hace que la máquina de estados sea comprensible, verificable y trazable con los requerimientos RF1/RF2.

\section{Conclusión}

Se definió un sistema de transición de estados con cinco estados, cuatro variables, condiciones de transición y actualizaciones de valores. El modelo se codificó en NuSMV y se verificaron propiedades de seguridad y vivacidad con \texttt{G} y \texttt{F}. La propiedad principal garantiza que toda solicitud válida eventualmente alcanza un resultado; las propiedades complementarias garantizan la coherencia entre el estado y el número de rutas, la terminalidad de la búsqueda y el retorno al estado inicial. El contraejemplo de la propiedad ``toda solicitud encuentra una ruta'' muestra que el modelo distingue correctamente entre una búsqueda sin resultados y una búsqueda exitosa.

\section{Cómo ejecutar}

\begin{enumerate}
  \item Descargar NuSMV desde \url{https://nusmv.fbk.eu/downloads.html}.
  \item Guardar el código en \texttt{rutas.smv}.
  \item Ejecutar \texttt{NuSMV rutas.smv}.
  \item Opcionalmente, abrir la consola interactiva de NuSMV y ejecutar \texttt{check\_ltlspec}.
\end{enumerate}

\section{Referencias}

\begin{itemize}
  \item NuSMV. \emph{NuSMV: A New Symbolic Model Checker}. \url{https://nusmv.fbk.eu/}
  \item Kochaleema, K. H.; Santhosh Kumar, G. \emph{Generic Methodology for Formal Verification of UML Models}. Defence Science Journal, 72(1), 40--48, 2022. DOI: 10.14429/dsj.72.17228.
  \item Hlaoui, Y. B.; Ayari, S. \emph{An Approach to Develop Correct-by-Construction Business Process Models Using a Formal Domain Specific Language}. Journal of Software: Evolution and Process, 37(2), e2762, 2025. DOI: 10.1002/smr.2762.
  \item Xavier, M.; Dubinin, V.; Patil, S.; Vyatkin, V. \emph{A Framework for the Generation of Monitor and Plant Model From Event Logs Using Process Mining for Formal Verification of Event-Driven Systems}. IEEE Open Journal of the Industrial Electronics Society, 5, 517--534, 2024. DOI: 10.1109/OJIES.2024.3406059.
  \item Nivon, Q.; Salaün, G. \emph{Debugging of BPMN Processes Using Coloring Techniques}. Lecture Notes in Computer Science, 2022, pp. 90--109. DOI: 10.1007/978-3-031-20872-0\_6.
\end{itemize}

\section{Anexo: modelo de incidencias}

Como ejemplo complementario, el modelo del ciclo de vida de una incidencia del foro se conserva en:
\begin{itemize}
  \item \texttt{Anexo\_Incidencias.md} y \texttt{Anexo\_Incidencias.tex}: informe del modelo de incidencias.
  \item \texttt{incidencias.smv}: código NuSMV del modelo de incidencias.
  \item \texttt{salida\_nusmv.txt}: salida de NuSMV para el modelo de incidencias.
\end{itemize}

Ese modelo no es el foco principal de este ejercicio; se mantiene porque el foro forma parte del proyecto y porque ilustra otro tipo de máquina de estados con variables de tiempo y notificación.

\end{document}
